\section{Introduction}
Image classifiers often learn shortcuts that work on most training images and fail on a minority of them \cite{geirhos2020shortcut}. Waterbirds \cite{sagawa2020groupdro} makes this concrete: birds from CUB \cite{wah2011cub} are pasted onto Places backgrounds \cite{zhou2018places} so that 95\% of training waterbirds appear on water and 95\% of landbirds on land. The standard metric is worst-group accuracy (WGA), the lowest accuracy over the four bird-by-background groups. The two smallest test groups contain 642 images each, so one image moves WGA by 0.16 percentage points (pp), and recent methods differ by a few tenths of a point.

A large share of recent progress comes from methods that keep an ERM backbone frozen and fit a light head on a small validation set with group labels \cite{kirichenko2023dfr,izmailov2022feature,labonte2023fewer}. Diverse prototypical ensembles (DPE) \cite{to2025dpe} are a strong recent example: fifteen prototype heads are trained in sequence on group-balanced subsets of the 1,199 annotated validation images, and their probabilities are averaged. For these methods the validation set is the whole of the supervision. Two things about it can limit them: how it is spent, and what it says.

\textbf{How it is spent.} Every extra step on top of DPE, such as learning weights for the heads, a combiner or a calibration map, needs labelled images of its own. Taking them from the 1,199 weakens the heads; reusing them makes the extra step optimistic. Cross-fitting \cite{chernozhukov2018dml} promises to avoid this trade-off, and calibration and threshold tuning \cite{platt1999probabilistic,elkan2001cost} are cheap. We test whether any of this beats DPE fitted on all images once the choice of setting is made on validation data only.

\textbf{What it says.} Three CUB species, Eastern Towhee, Western Meadowlark and Western Wood-Pewee, are labelled waterbird in the Waterbirds metadata although they are landbirds. The generating script marks a species as waterbird if its name contains a keyword such as ``tern'', and this string also occurs in ``Eastern'' and ``Western''. The error is known: the dataset repository lists the three species and calls it harmless for the validity of the benchmark, and MaskTune \cite{asgari2022masktune} corrected it. Its effect on methods that are \emph{trained} on validation group labels has not been measured. For DPE it is not small: 11 of the 133 validation images of the waterbird-on-land group, the group DPE fails on most, carry the wrong label.

We study both questions on seven saved ResNet-50 backbones (five standard ERM, two longer-trained ERM$^*$) and report three findings.
\begin{itemize}
\item \textbf{The label error dominates DPE's worst group.} These species cause 50--62\% of DPE's errors on waterbird-on-land. Refitting DPE after correcting the 18 affected validation images raises WGA on corrected test labels by $+2.48$\pp{} (95\% interval $[+1.80,+3.15]$, all five standard-ERM backbones) and by about $+2.1$\pp{} on ERM$^*$. The official labels score the same fix as $-0.83$\pp: a benchmark that keeps the error penalises the correction.
\item \textbf{Spending the validation set more carefully helps little.} With validation-only selection repeated over ten split seeds, cross-fitting with recalibration is $-0.17$\pp{} ($[-0.33,-0.02]$) relative to full-validation DPE, learned combiners on a held-out split are worse, and recalibration is equivalent to a near-zero threshold shift. A $+0.47$\pp{} gain that we had highlighted after looking at test results shrinks to $+0.17$\pp{} over ten split seeds. Head-seed noise alone is about 0.4\pp.
\item \textbf{Label errors change method comparisons.} Correcting the test labels turns a $+0.54$\pp{} advantage of a cross-fitted ensemble over DPE into $-0.18$\pp. With the official labels, DPE is 1.5\pp{} ahead of DFR-style retraining on the same features and labels; both gain from cleaner validation labels.
\end{itemize}
We also reproduce DPE with the authors' code on two fresh ERM backbones (89.4\% and 86.5\% WGA against 91.0\% reported) and release all code, per-run results and audits.\footnote{\url{https://github.com/chautanphat/waterbirds-shortcut-mitigation}}

\textbf{Scope.} We do not propose a new method or claim a state of the art. The study covers one dataset, the same backbones and test set were inspected during earlier development, and the label correction concerns three species only. The practical message is narrow but direct: for validation-label methods on Waterbirds, cleaning 18 labels matters more than any of the aggregation schemes we tried, and comparisons should report corrected labels alongside the official ones.
